\section{Benchmark Study: LLM vs.\ VLM Task Planning}
\label{sec:benchmark}

\subsection{Motivation and design}

With both a static-scene LLM pipeline and a camera-driven VLM pipeline
working end-to-end (Section~\ref{sec:architecture}), this study ran a direct
head-to-head comparison on real hardware. The design is adapted from the
task-length / goal-specificity / affordance-perception taxonomy and the
TS/TSR/AETS metrics (Eq.\ 14--16) of Favali, Sabattini, and Villani
(RO-MAN 2025), scoped down to what is runnable on this physical UR5cb setup
with one LLM model, one VLM model, and no dynamic-event injection. Full
methodology lives in \texttt{benchmark/PLAN.md}; this section summarizes it
and reports the results.

\paragraph{What differs from the source paper.} The paper validates in a
symbolic simulator where ground-truth world state is always exactly known,
so both safety and success metrics are fully automatic. On real hardware,
neither is: there is no collision/force-fault sensor anywhere in
\texttt{ur5\_skill\_executor\_node.py}, and there is no vision-based
final-state verification step. The benchmark is therefore
\textbf{semi-automatic}: everything the program can know (command, model,
config, generated plan, steps actually executed, service-level errors) is
logged automatically; the operator supplies exactly two numbers per trial,
immediately after execution -- \emph{was it safe?} and \emph{how many
sub-tasks actually completed correctly?} -- via an inline GUI form or the
standalone \texttt{benchmark/benchmark\_annotate.py} script, which computes
TS/TSR/AETS from those two numbers plus the automatically-logged step count.

\paragraph{Rerun in September, not an extension of the July study.} An
earlier version of this benchmark ran 15--20 July 2026 on
\texttt{nebius/nvidia-nemotron-120b} (LLM) and \texttt{groq/qwen3.6-27b}
(VLM). By September, \texttt{qwen3.6-27b} had left Nebius's model catalog
with no same-scale replacement (ROBOAI-12); rather than mix a new VLM model
into a comparison whose LLM side stayed on the old one, both arms were
rerun from scratch on a single unified model. This is a new baseline, not
directly comparable to the July numbers point-for-point -- see the
Findings subsection below for a case where the two studies actually
disagree, not just differ in magnitude. A
third arm, \texttt{VLM\_LLM} (VLM scene-grounding feeding an LLM planner),
was attempted alongside this rerun but abandoned after 3 trials when a
real object in the camera's field of view (a loose cardboard flap propped
on a tray-like container, outside the intended workspace) returned no
usable depth at its viewing angle and crashed the hybrid grounding call
every time the VLM detected it -- a physical-setup problem, not a code
defect, and out of scope to chase mid-campaign. Those 3 trials
(\texttt{pp\_easy}, \texttt{pp\_hard}, \texttt{sort\_easy}) remain in
\texttt{benchmark/results\_2026-09\_3arm.csv} as partial data but are
excluded from every table and figure below.

\subsection{Task matrix}

Six conditions across three difficulty axes (task length, reference
specificity, affordance/perception load), each at an ``easy'' and ``hard''
extreme, over the real scene: four identically-shaped cubes distinguished
only by color (blue, red, white, orange) and a tray as the placement target.

\begin{longtable}{@{}p{1.6cm}p{6.7cm}p{1.5cm}p{1.6cm}p{1.9cm}@{}}
\toprule
\textbf{Task} & \textbf{Prompt} & \textbf{Length} & \textbf{Specificity} & \textbf{Sub-tasks} \\
\midrule
\endhead
\texttt{pp\_easy} & ``Pick the red cube and place it on the brown tray.'' &
  short (4 steps) & specific & 1 \\
\texttt{pp\_hard} & ``Put all the cubes on the brown tray.'' &
  long (16 steps) & lifted & 4 \\
\texttt{sort\_easy} & ``Stack the red cube on top of the blue cube.'' &
  short (4 steps) & specific & 1 \\
\texttt{sort\_hard} & ``Sort all the cubes: put the red cube onto the brown
  tray, and arrange the rest in a straight line on the table.'' &
  long (16 steps) & lifted & 4 \\
\texttt{arith\_easy} & ``Each cube has a value: blue = 1, red = 2, white = 3,
  orange = 4. Pick up the cube whose value equals 3 minus 1, and place it on
  the brown tray.'' & short (4 steps) & specific (computed) & 1 \\
\texttt{arith\_hard} & ``Each cube has a value: blue = 1, red = 2, white = 3,
  orange = 4. Move every cube whose value is greater than the sum of the
  blue and red cubes' values onto the brown tray, and arrange the rest in a
  line on the table.'' & long (16 steps) & lifted (computed) & 4 \\
\bottomrule
\end{longtable}

\noindent The arithmetic task exists specifically to test numeric reasoning
grounded in the scene rather than pure spatial pick-place: \texttt{arith\_easy}
needs one subtraction to resolve a single unambiguous cube; \texttt{arith\_hard}
needs an addition to compute a threshold, then a per-cube comparison against
it -- genuinely ``lifted'' (no cube is named directly) and requiring the
model to work out its own two-group split.

\textbf{Design}: 6 conditions $\times$ 10 repetitions $\times$ 2 models
(LLM and VLM, both \texttt{nebius/kimi-k2.6}) = \textbf{120 trials}. A
single reasoning method, \texttt{cot\_sc}, was held fixed across all trials
-- per the source paper it is the most safety-robust method, so it avoids
stacking reasoning-method variance on top of the LLM-vs-VLM comparison
actually being measured. Trials were collected 4--9 September 2026, on
branch \texttt{feature/ROBOAI-19-bbox-scene-schema} (commit history in
ROBOAI-25); a code freeze was in effect for the duration of collection, with
one logging-only exception (ROBOAI-29, a stale-CSV-header bug with no effect
on planner/executor behavior).

\subsection{Metrics}

\begin{itemize}[nosep]
  \item \textbf{TS} (Task Safety): 1 if the trial was safe (no
    collision/damage), 0 otherwise -- operator-supplied.
  \item \textbf{TSR} (Task Success Rate):
    $\text{sub\_tasks\_completed} / \text{sub\_tasks\_required}$ --
    \texttt{sub\_tasks\_required} is fixed per task (table above),
    \texttt{sub\_tasks\_completed} is operator-supplied after inspecting the
    final scene.
  \item \textbf{AETS} (Action Efficiency on Task Success):
    $\text{sub\_tasks\_completed} / (\text{sub\_tasks\_required} \times
    \text{steps\_executed})$ -- fully automatic given the two manual inputs
    above, since \texttt{steps\_executed} is read from the logged execution
    result.
\end{itemize}

At $n=10$ per cell, point estimates alone invite reviewer pushback (per
ROBOAI-28). Every reported TSR is also given as a Wilson score interval at
95\% confidence (pooled at the sub-task level, i.e.\ successes over
$\text{sub\_tasks\_required} \times n$, not over trial count -- consistent
with AETS's own denominator). One comparison was planned before collection,
not chosen afterward to fit the numbers: a two-sided Fisher's exact test,
LLM vs.\ VLM, on \texttt{arith\_hard} TSR, at $\alpha=0.05$. Everything else
is reported descriptively only, per \texttt{benchmark/PLAN.md}~\S3.

\subsection{Results}

\begin{table}[H]
\centering
\caption{Overall results, 60 trials per model (120 total).}
\label{tab:bench-overall}
\begin{tabular}{@{}lcccc@{}}
\toprule
\textbf{Model} & \textbf{n} & \textbf{TS (safety)} & \textbf{TSR (success)} & \textbf{AETS (efficiency)} \\
\midrule
LLM (\texttt{kimi-k2.6}) & 60 & 100\% & 98.3\% & 0.1269 \\
VLM (\texttt{kimi-k2.6}) & 60 & 100\% & 89.2\% & 0.1149 \\
\bottomrule
\end{tabular}
\end{table}

\begin{table}[H]
\centering
\caption{Results by task and model. Steps = mean steps executed.}
\label{tab:bench-by-task}
\begin{tabular}{@{}llccccc@{}}
\toprule
\textbf{Task} & \textbf{Model} & \textbf{n} & \textbf{TS} & \textbf{TSR} & \textbf{AETS} & \textbf{Steps} \\
\midrule
\texttt{pp\_easy}    & LLM & 10 & 100\% & 100\% & 0.2000 & 5.0 \\
\texttt{pp\_easy}    & VLM & 10 & 100\% & 100\% & 0.2000 & 5.0 \\
\texttt{pp\_hard}    & LLM & 10 & 100\% & 100\% & 0.0579 & 17.3 \\
\texttt{pp\_hard}    & VLM & 10 & 100\% &  93\% & 0.0526 & 17.6 \\
\texttt{sort\_easy}  & LLM & 10 & 100\% & 100\% & 0.2000 & 5.0 \\
\texttt{sort\_easy}  & VLM & 10 & 100\% &  70\% & 0.1400 & 5.0 \\
\texttt{sort\_hard}  & LLM & 10 & 100\% & 100\% & 0.0558 & 18.0 \\
\texttt{sort\_hard}  & VLM & 10 & 100\% &  88\% & 0.0505 & 17.4 \\
\texttt{arith\_easy} & LLM & 10 & 100\% & 100\% & 0.1967 & 5.1 \\
\texttt{arith\_easy} & VLM & 10 & 100\% & 100\% & 0.2000 & 5.0 \\
\texttt{arith\_hard} & LLM & 10 & 100\% &  90\% & 0.0509 & 17.7 \\
\texttt{arith\_hard} & VLM & 10 & 100\% &  85\% & 0.0461 & 18.6 \\
\bottomrule
\end{tabular}
\end{table}

\begin{table}[H]
\centering
\caption{Wilson 95\% CI (pooled sub-tasks) and the one pre-registered test.}
\label{tab:bench-wilson}
\begin{tabular}{@{}llcc@{}}
\toprule
\textbf{Task} & \textbf{Model} & \textbf{TSR} & \textbf{95\% CI} \\
\midrule
\texttt{arith\_hard} & LLM & 90.0\% & [76.9, 96.0] \\
\texttt{arith\_hard} & VLM & 85.0\% & [70.9, 92.9] \\
\midrule
\multicolumn{4}{@{}l}{Fisher's exact test, LLM vs.\ VLM, \texttt{arith\_hard} TSR: $p = 0.74$.} \\
\multicolumn{4}{@{}l}{Overall pooled TSR: LLM 97.3\% [93.3, 99.0], VLM 88.7\% [82.6, 92.8].} \\
\bottomrule
\end{tabular}
\end{table}

\begin{figure}[H]
\centering
\includegraphics[width=0.85\textwidth]{overall_summary.png}
\caption{Overall TS/TSR/AETS summary, LLM vs.\ VLM.}
\label{fig:overall-summary}
\end{figure}

\begin{figure}[H]
\centering
\includegraphics[width=0.85\textwidth]{ts_tsr_by_task.png}
\caption{Task Safety and Task Success Rate broken down by task condition.}
\label{fig:ts-tsr-by-task}
\end{figure}

\begin{figure}[H]
\centering
\includegraphics[width=0.85\textwidth]{easy_vs_hard.png}
\caption{Easy vs.\ hard difficulty comparison, aggregated across task families.}
\label{fig:easy-vs-hard}
\end{figure}

\begin{figure}[H]
\centering
\includegraphics[width=0.85\textwidth]{aets_by_task.png}
\caption{Action efficiency (AETS) by task condition.}
\label{fig:aets-by-task}
\end{figure}

\begin{figure}[H]
\centering
\includegraphics[width=0.85\textwidth]{steps_by_task.png}
\caption{Mean executed plan length by task condition.}
\label{fig:steps-by-task}
\end{figure}

\begin{figure}[H]
\centering
\includegraphics[width=0.85\textwidth]{issue_rate_by_task.png}
\caption{Rate of logged operator issue notes by task condition.}
\label{fig:issue-rate-by-task}
\end{figure}

\begin{figure}[H]
\centering
\includegraphics[width=0.7\textwidth]{failure_modes.png}
\caption{Failure-mode breakdown, derived from the 22 non-empty operator
  notes out of 120 trials.}
\label{fig:failure-modes}
\end{figure}

\subsection{Findings}

\paragraph{Headline.} The LLM pipeline is both safer and more successful
than the VLM pipeline overall (100\%/98.3\% vs.\ 100\%/89.2\%), but the gap
is real and modest, not dramatic -- both arms are safe every single trial,
and the pre-registered comparison on \texttt{arith\_hard} (the condition
the July study flagged as the sharpest divide) shows no detectable
difference ($p=0.74$).

\paragraph{The dominant failure mode is not vision, and it is not
value-mapping arithmetic either -- it is multi-object placement geometry,
in both arms.} Reading all 22 non-empty operator notes (rather than
trusting the aggregate numbers alone) was the key analysis step here, as it
was in July, but the finding is different this time. Every one of the 22
notes describes a collision between objects or a mispositioned/dropped
release during a multi-cube task (\texttt{pp\_hard}, \texttt{sort\_hard},
\texttt{arith\_hard}) -- ten collisions, twelve grasp/execution failures,
\textbf{zero} notes describing a wrong object chosen for reasoning reasons,
and zero incomplete-enumeration failures (Figure~\ref{fig:failure-modes}).
This happens in \textbf{both} arms: LLM contributes 4 of the 22 notes (all
on \texttt{arith\_hard}), VLM the rest, spread across all three
multi-object conditions roughly in proportion to their trial counts. The
failure class recurs regardless of which planner grounded the scene,
consistent with a genuine geometry-grounding gap (specifically: nothing in
the pipeline prevents the arm from clipping an already-placed object while
approaching or releasing the next one, even when the two objects' final
target positions do not overlap) rather than a VLM-specific weakness.

\paragraph{This does not replicate July's arithmetic-collapse finding.} The
July study (different VLM model, \texttt{qwen3.6-27b}) found VLM TSR
collapsing to 30\%/25\% on \texttt{arith\_easy}/\texttt{arith\_hard} against
100\%/97\% for LLM, and traced 12 of 13 ``wrong cube'' notes to an
apparent off-by-one in the stated value mapping. On \texttt{kimi-k2.6},
\texttt{arith\_easy} is 100\%/100\% and \texttt{arith\_hard} is 90\%/85\%
(Table~\ref{tab:bench-by-task}), and there are no ``wrong cube'' notes at
all in September's data. The simplest reading is that July's finding was a
property of \texttt{qwen3.6-27b} specifically, not of VLM grounding in
general -- consistent with this project's other observed instances of
model-specific behavioral drift (Section~\ref{sec:tech-decisions}). This
matters for how the two studies' results should be used together: July's
data is not wrong, but its central claim does not generalize across models,
which is itself worth stating plainly rather than quietly dropping.

\paragraph{Plan length is not the differentiator.} LLM and VLM produce plans
within about one step of each other for every task condition
(Table~\ref{tab:bench-by-task}, Figure~\ref{fig:steps-by-task}) -- the gap is
in what happens physically after the plan is formed, not in the plan
itself.

\subsection{Reproducing or extending this study}

\begin{itemize}[nosep]
  \item Raw data: \texttt{benchmark/results\_2026-09\_3arm.csv} (one row per
    trial; includes 3 leftover \texttt{VLM\_LLM} rows, excluded from every
    table/figure above -- see the design note).
  \item Plan and full methodology: \texttt{benchmark/PLAN.md}.
  \item Figure generation: \texttt{benchmark/plot\_results.py} (regenerates
    all PNGs in \texttt{benchmark/figures/} from the results CSV).
  \item Annotation tooling: \texttt{benchmark/benchmark\_annotate.py}
    (standalone CLI) or the in-GUI ``Benchmark trial'' checkbox and inline
    form (\texttt{bridge\_node.py::record\_benchmark\_annotation}).
\end{itemize}

The most direct extension -- attempted this cycle, not completed -- is the
\texttt{VLM\_LLM} hybrid third arm; see the design note above for why it was
abandoned after 3 trials. A validation ablation (deterministic
release-geometry fixes on vs.\ off, on the three ``hard'' conditions, where
this section's dominant failure mode actually lives) is planned but not yet
run -- see ROBOAI-30.
